% applications/vehicle_space_and_platforms.tex

\chapter{Vehicle Space and Platform Interaction}

\label{chap:applications-vehicle-space-platform}

This chapter applies canonical pose3 and runtime interpretation to
vehicle-space and platform-edge use cases.

% ============================================================
\section{Application Context}
% ============================================================

Platform interaction and clearance assessment depend on vehicle-space
behaviour relative to railway runtime state. The relevant engineering
object is therefore not centerline geometry alone, but the derived
vehicle-relative state evaluated from canonical runtime services.

% ============================================================
\section{Derived Vehicle-Space Interpretation}
% ============================================================

Vehicle-space behaviour is interpreted downstream of canonical pose3 and
frame/cant evaluation. It is application interpretation, not a new
state-model layer.

This enables context-specific evaluation of platform gaps, clearance,
and envelope behaviour under operational conditions.

% ============================================================
\section{Engineering Use Cases}
% ============================================================

Typical use cases include platform-edge verification, envelope checks in
curves, and scenario-based assessment for differing vehicle classes.
These use cases consume canonical runtime and realization layers.

% ============================================================
\section{Representation Boundary}
% ============================================================

Rendered or exported geometry is treated as downstream representation of
runtime-evaluated state, not as engineering truth by itself.

% ============================================================
\section{Connection to AIM}
% ============================================================

\begin{summary}
Vehicle-space and platform applications in AIM are downstream
interpretations built on canonical pose3 and runtime services.

They demonstrate practical decision support while preserving the
separation between engineering objects and representations.
\end{summary}
